% Original1945 Hempel treatment; Sprenger2011 is a distinct unheld allocation.
\chapter{Hempel's Confirmation Conditions and Their Paradoxes}
\label{sec:p15-hempel-confirmation-and-paradoxes}

Hempel asks how a stated body of evidence bears on a stated hypothesis.
He does not offer rules which mechanically invent a scientific theory from
observations, nor identify logical confirmation with a feeling of conviction.
His 1945 study separates a nonquantitative relation of confirmation from
degrees of confirmation, the reliability of an observation report, and the
decision to accept a hypothesis. Those are different questions even when
one investigation raises all of them
\parencite[I, \S\S1--2, pp.~1--9]{Hempel1945ConfirmationI}.

The sources here are the two original \emph{Mind} articles, printed
pp.~1--26 and 97--121. Their definitions, arguments and qualifications are
not substituted for Jan Sprenger's complete 2011 historical chapter,
pp.~235--263, which is a distinct historical analysis rather than either
original essay
\parencite{Sprenger2011HempelConfirmation}. The nearby Goodman treatment
addresses projectibility; the probability comparisons in
\cref{sec:p15-r43-rankings} address a later, explicitly modeled notion of
incremental support. Neither answers Hempel's question merely by sharing
the word \emph{confirmation}.

\section{Reports, hypotheses and the kind of conclusion sought}
\label{sec:p15-hempel-task}

\begin{hbdefinition}{def:p15-hempel-reports}{The stated evidence and the observational language}
Fix a first-order relational language. Its predicates belong to a specified
observational vocabulary; it has individual names, Boolean connectives and
individual quantifiers, but no functions, identity or predicate variables.
An observation report is initially a finite conjunction of positive or
negative ground atomic sentences. Hempel also considers the broader class
of quantifier-free Boolean reports. A hypothesis may have any quantifier
structure expressible in this language.

Write $E$ for a report, $H$ for a hypothesis, and $C(E,H)$ for confirmation.
Hempel uses $B$ for the report in Part~II; this is only a change of report
metavariable. The relation is evaluated on the \emph{total stated report}.
Observational vocabulary is relative to the agreed observation techniques,
not necessarily to unaided perception. Being an observation-type sentence
does not imply being true, actually observed, or reliable.
\end{hbdefinition}

This logical comparison can describe what evidence would bear on $H$ if it
were available. Scientific acceptance additionally requires a reliable
record and a policy of acceptance. Hempel leaves the pragmatic reliability
question outside his formal construction. He also distinguishes a logical
relationship between sentences from the experimental activity which leads
to accepting one of those sentences
\parencite[I, \S6, pp.~22--26; II, \S10, pp.~114--119]{Hempel1945ConfirmationI,Hempel1945ConfirmationII}.

\section{Why the instance rule depends on a spelling}
\label{sec:p15-hempel-nicod}

\begin{hbdefinition}{def:p15-hempel-nicod}{Nicod's displayed criterion and unary sufficiency}
For a displayed universal conditional $\forall x(P(x)\to Q(x))$, Nicod's
criterion classifies a reported $P(a)\land Q(a)$ as confirming,
$P(a)\land\neg Q(a)$ as disconfirming, and a case with $\neg P(a)$ as
neutral. As a necessary-and-sufficient criterion this inspects the
antecedent and consequent of a particular formulation.

The weaker \emph{unary instance-sufficiency stipulation} retains only that
$P(a)\land Q(a)$ confirms the displayed unary conditional. Hempel explicitly
restricts this stipulation to one universally quantified variable; it is
not an unrestricted rule about any tuple satisfying a conditional matrix.
These are proposed standards for a relation, not conclusions proved by
deductive logic alone.
\end{hbdefinition}

\begin{hbtheorem}{thm:p15-hempel-nicod-spelling}{proved=annex:p15-hempel-nicod-spelling}{Equivalent raven sentences and the four cases}
Let $\mathsf R$ mean raven and $\mathsf B$ mean black. The sentences
\[
 H_R=\forall x(\mathsf R(x)\to\mathsf B(x)),\qquad
 H_N=\forall x(\neg\mathsf B(x)\to\neg\mathsf R(x))
\]
are logically equivalent. For the four complete specimen reports, the
displayed Nicod classification is:
\[
\begin{array}{c|cc}
 \text{reported properties}&H_R&H_N\\\hline
 \mathsf R(a),\mathsf B(a)&\text{confirming}&\text{neutral}\\
 \mathsf R(a),\neg\mathsf B(a)&\text{disconfirming}&\text{disconfirming}\\
 \neg\mathsf R(a),\mathsf B(a)&\text{neutral}&\text{neutral}\\
 \neg\mathsf R(a),\neg\mathsf B(a)&\text{neutral}&\text{confirming}
\end{array}
\]
The equivalent formulation
$\forall x((\mathsf R(x)\land\neg\mathsf B(x))\to
(\mathsf R(x)\land\neg\mathsf R(x)))$ has no Nicod-positive case.
Another equivalent formulation is
\[
 \forall x((\mathsf R(x)\lor\neg\mathsf R(x))\to
                       (\neg\mathsf R(x)\lor\mathsf B(x))).
\]
Unary instance sufficiency applied to this last spelling, together with
equivalence of confirmed hypotheses, counts every reported nonraven and
every reported black object as confirming $H_R$.
\end{hbtheorem}

The inference is conditional on two proposed confirmation principles. It
is not a proof that observing a red pencil makes the raven law probable
to a particular numerical degree. Conversely, rejecting Nicod's neutrality
clause does not require rejecting its restricted sufficiency clause.
Hempel develops precisely this distinction
\parencite[I, \S\S3--5, pp.~9--15]{Hempel1945ConfirmationI}.

\begin{hbtheorem}{thm:p15-hempel-multiple-instance}{proved=annex:p15-hempel-multiple-instance}{A positive ordered pair can refute the universal sentence}
Put
\[
 H_*=\forall x\forall y\bigl(
 \neg(R(x,y)\land R(y,x))\to(R(x,y)\land\neg R(y,x))\bigr).
\]
Then $H_*$ is equivalent to $\forall x\forall y R(x,y)$.
On the two-element domain $\{a,b\}$, with $R=\{(a,b)\}$, the pair
$(a,b)$ satisfies both the displayed antecedent and consequent, while the
report $R(a,b)\land\neg R(b,a)$ entails $\neg H_*$. Thus unrestricted
tuple-instance sufficiency conflicts with the requirement that a consistent
confirming report be compatible with the confirmed hypothesis.
\end{hbtheorem}

This is Hempel's own qualification, not a later probability objection.
It prevents extending the unary raven argument indiscriminately to
relations of arbitrary arity
\parencite[I, p.~13, n.~1]{Hempel1945ConfirmationI}.

\section{The proposed repairs and Hempel's actual reply}
\label{sec:p15-hempel-raven-responses}

\begin{hbtheorem}{thm:p15-hempel-existential-import}{proved=annex:p15-hempel-existential-import}{Adding existential import changes the hypothesis}
Although $H_R$ and $H_N$ are equivalent, the sentences
\[
 H_R\land\exists x\mathsf R(x),\qquad
 H_N\land\exists x\neg\mathsf B(x)
\]
are not. A singleton black-raven model satisfies the first but not the
second; a singleton nonblack-nonraven model satisfies the second but not
the first.
\end{hbtheorem}

Existential import can clarify a speaker's intended claim, but it is not
a general repair preserving the original universal law. Hempel notes
that a test hypothesis may leave ambiguous whether the existential clause
concerns persons, persons injected with a substance, or persons reacting
positively. Hypotheses about unrealized extreme conditions need not assert
that their antecedents occur at all. His examples are arguments about
interpretation, not claims that the described medical or biological tests
are empirically correct
\parencite[I, \S5.11, pp.~15--17]{Hempel1945ConfirmationI}.

A second proposal adds a privileged field of application, such as the class
of ravens. Hempel objects that this extra parameter is not uniquely read
off from the universal sentence. A sodium-salt law is also used negatively:
if a specimen fails the yellow-flame test, the law excludes sodium salt.
Restricting attention to already identified sodium salts loses that use.
Adding a field can define a different inferential task; merely naming it
does not show that it preserves the task originally under discussion
\parencite[I, \S5.12, pp.~17--18]{Hempel1945ConfirmationI}.

Hempel's preferred reply keeps the universal conditional. Such a sentence
forbids a $P$ which is not a $Q$ anywhere in its domain; it is not silent
about every non-$P$. He further contrasts an unknown substance, whose
negative flame test is followed by a finding that it contains no sodium
salt, with ice already known to contain none. In the second experiment,
either flame outcome is compatible with the prior non-sodium information.
One must not compare $H$ with $E$ alone while tacitly importing that prior
information. His characterization of the appearance of paradox as a
psychological illusion is an attributed philosophical diagnosis, not a
newly proved psychological law
\parencite[I, \S5.2, pp.~18--21]{Hempel1945ConfirmationI}.

The footnote on Hosiasson-Lindenbaum discusses a different question: how
large the support increment is. Hempel notes that an assumed numerical
imbalance between ravens and nonblack objects is not fixed by logic, and
can behave differently when individuals are space-time regions rather than
ordinary things. The later finite probability model in
\cref{sec:p15-r43-rankings} declares its marginals and observation model;
it must not be retrospectively made Hempel's 1945 definition
\parencite[I, pp.~21--22, n.~2]{Hempel1945ConfirmationI}.

\section{Successful prediction does not define all confirmation}
\label{sec:p15-hempel-prediction}

\begin{hbdefinition}{def:p15-hempel-prediction}{The prediction criterion}
Represent a report by a finite set of observation sentences. The prediction
criterion calls it confirming when it can be partitioned into $E_1,E_2$,
with $E_2$ nonempty, so that every sentence of $E_2$ follows from
$E_1\cup\{H\}$ but not from $E_1$ alone. It calls the report disconfirming
when the report and $H$ are inconsistent. This is the source's candidate
criterion, not the satisfaction definition introduced below.
\end{hbdefinition}

\begin{hbtheorem}{thm:p15-hempel-prediction-gap}{proved=annex:p15-hempel-prediction-gap}{Two defects of prediction as a general criterion}
In the relational language, let
\[
 J=\forall x((\forall y R_1(x,y))\to\exists zR_2(x,z)).
\]
For a finite ground observation report $E$ and any ground quantifier-free
sentence $A$, $E\land J\models A$ implies $E\models A$, when the admissible
domains are infinite. Thus $J$ supplies no new ground observation
consequences from such a report.

Separately, take the two-sentence report $\{P(a),Q(a)\}$ and
$H=P(a)\to Q(a)$. It satisfies the prediction criterion for $H$, but
not for the weaker hypothesis $H\lor T(c)$, where $T(c)$ is an independent
ground atom. The criterion consequently does not preserve confirmation
under passage to weaker hypotheses.
\end{hbtheorem}

Hempel's first point concerns an antecedent which itself has general
content. A few recorded $R_1$ cases may support accepting $\forall yR_1(a,y)$
without entailing it; a predicted existential $R_2$ claim may in turn be
supported by a named witness. His \emph{quasi-induction} is acceptance of
an already given intermediate hypothesis, not its mechanical discovery.
Replacing deduction in the prediction criterion by deduction plus these
steps would presuppose the confirmation relation being defined. His
polarized-light example similarly points to theoretical and auxiliary
claims between observations and prediction
\parencite[II, \S7, pp.~97--102]{Hempel1945ConfirmationII}.

\section{Conditions of adequacy are stipulated constraints}
\label{sec:p15-hempel-conditions}

\begin{hbdefinition}{def:p15-hempel-conditions}{Entailment, consequence, equivalence and consistency}
Hempel's principal standards for $C$ are:
\begin{enumerate}
\item \emph{Entailment}: if $E\models H$, then $C(E,H)$.
\item \emph{General consequence}: if $C(E,K)$ holds for every $K$ in a
      class $\Gamma$, and $\Gamma\models H$, then $C(E,H)$.
\item \emph{Consistency}: for consistent $E$, the class of all hypotheses
      confirmed by $E$ is jointly compatible with $E$.
\end{enumerate}
General consequence includes confirmation of a consequence of one confirmed
hypothesis, conjunction of confirmed hypotheses, and substitution of a
logically equivalent hypothesis. These are specifications by which a
proposal is to be judged, not deductions that every scientifically useful
notion of evidence must satisfy them. Contradictory reports are excluded
from the consistency clause: entailment would make them confirm everything.
\end{hbdefinition}

\begin{hbtheorem}{thm:p15-hempel-converse-collapse}{proved=annex:p15-hempel-converse-collapse}{Unrestricted upward and downward closure collapse}
Suppose $C$ satisfies entailment and special consequence, and also
\emph{converse consequence}: if $C(E,H)$ and $K\models H$, then $C(E,K)$.
Then $C(E,H)$ holds for every report $E$ and every hypothesis $H$.
\end{hbtheorem}

Hempel's argument uses $E$, then $E\land H$, then $H$. It explains why
successful predictions cannot license confirmation of \emph{every} stronger
theory. The empirical role of a law's specialization, possibly with auxiliary
assumptions, is not the arbitrary addition of any conjunct. This is the
three-condition argument in the original article; the neighboring
\cref{thm:p15-goodman-confirmation-collapse} packages different initial
closure assumptions
\parencite[II, \S8, pp.~102--105]{Hempel1945ConfirmationII}.

\begin{hbtheorem}{thm:p15-hempel-joint-consistency}{proved=annex:p15-hempel-joint-consistency}{Individual compatibility cannot liberalize unrestricted joint closure}
Assume general consequence and require, for consistent $E$, that each
confirmed $H$ be individually compatible with $E$. These assumptions
already imply that all hypotheses confirmed by $E$ are jointly compatible
with it.

By contrast, individual fit to a finite record need not imply joint
compatibility of proposed laws. The functions $f_0(x)=0$ and
$f_1(x)=x(x-1)$ both fit the observations $(0,0),(1,0)$, but disagree at
$x=2$ and cannot both be one function on $\{0,1,2\}$.
\end{hbtheorem}

Hempel raises the curve-fitting objection and considers retaining only
individual compatibility while abandoning general consistency. Under the
unrestricted general-consequence condition stated above, this is not an
independent weakening. To permit mutually incompatible fitted hypotheses
to be confirmed by the same report one must also restrict joint consequence
closure. This logical qualification leaves his philosophical choice between
stricter and more liberal reconstructions visible; it does not ascribe a
new theorem to his proposed liberalization. He also observes that bare
entailment meets the formal conditions while being too narrow for the
intended material conception of confirmation
\parencite[II, pp.~105--107]{Hempel1945ConfirmationII}.

\section{Development, direct confirmation and its consequence closure}
\label{sec:p15-hempel-satisfaction}

\begin{hbdefinition}{def:p15-hempel-development}{Development on a finite named class}
Let $D$ be a finite nonempty set of individual names. Define
$\operatorname{dev}_D(H)$ recursively: leave ground atoms unchanged,
commute with Boolean operations, replace $\forall xA$ by
$\bigwedge_{d\in D}\operatorname{dev}_D(A[d/x])$, and replace $\exists xA$
by $\bigvee_{d\in D}\operatorname{dev}_D(A[d/x])$. Substitution is
capture-avoiding. The result states what $H$ says when the quantified
individuals range over the named class. It is not the assertion that the
actual scientific domain contains only that class.
\end{hbdefinition}

\begin{hblemma}{thm:p15-hempel-development}{proved=annex:p15-hempel-development}{The finite-domain meaning of development}
In a nonempty structure whose domain is covered by the denotations of $D$,
every sentence $H$ is equivalent to $\operatorname{dev}_D(H)$.
Distinct names need not denote distinct objects for this equivalence.
Outside that domain-covering hypothesis, agreement need not hold.
\end{hblemma}

For example, the development of $\forall x\exists y L(x,y)$ on $\{a,b\}$ is
\[
 (L(a,a)\lor L(a,b))\land(L(b,a)\lor L(b,b)).
\]
The development of $\exists xP(x)$ is $P(a)\lor P(b)$; that of a
quantifier-free $P(c)\lor Q(c)$ remains the same sentence, even when
$c\notin D$. Hempel uses these cases to extend the method beyond simple
universal conditionals
\parencite[II, \S9, pp.~108--110]{Hempel1945ConfirmationII}.

\begin{hbdefinition}{def:p15-hempel-essential}{Essential mention and the declared fragment}
Treat the finitely many ground atoms of a quantifier-free report as
propositional inputs. An atom is \emph{essential} when changing its truth
value can change the report's Boolean value. Equivalently, it cannot be
eliminated from every equivalent Boolean formulation. Let $D(E)$ contain
the individual names occurring in essential atoms.

For the explicit construction below, $E$ is consistent and $D(E)$ is
nonempty. There is no identity, no function symbol, no background
restriction identifying names, and no additional theory restricting
predicate interpretations. These are substantive assumptions.
Purely analytic reports with empty essential support and arbitrary richer
languages are outside the following construction; no empty scientific
domain is silently introduced.
\end{hbdefinition}

\begin{hbdefinition}{def:p15-hempel-satisfaction-core}{The specified satisfaction construction}
Put
\[
 \Delta_E=\{K:E\models\operatorname{dev}_{D(E)}(K)\},\qquad
 C(E,H)\quad\Longleftrightarrow\quad\Delta_E\models H.
\]
Membership in $\Delta_E$ is \emph{direct confirmation}; its consequence
closure gives confirmation. Disconfirmation of $H$ means confirmation
of $\neg H$; neutrality means neither. These are semantic definitions,
not an asserted terminating decision algorithm for arbitrary hypotheses.
They implement the nonempty named-evidence fragment of Hempel's corrected
(9.1a) and (9.2), with the stated assumptions made explicit.
\end{hbdefinition}

\begin{hbtheorem}{thm:p15-hempel-essential-repair}{proved=annex:p15-hempel-essential-repair}{A tautological added name and the corrected report class}
The reports $E=P(a)$ and $E'=P(a)\land(Q(b)\lor\neg Q(b))$ are equivalent.
If every syntactically mentioned name is used, $E$ directly confirms
$\forall xP(x)$ but $E'$ does not; under that uncorrected consequence
closure $E'$ does not confirm it either. Using essential mention gives
$D(E)=D(E')=\{a\}$ and the same confirmation classification.
More generally, equivalent admissible reports have the same essential
ground atoms and the same $D(E)$.
\end{hbtheorem}

The correction is already in Part~II, not a later repair invented here.
Hempel credits a discussion with Olaf Helmer and distinguishes conjunctions
of observation literals from the broader Boolean-report convention used
in his earlier technical article
\parencite[II, pp.~110--111, n.~1]{Hempel1945ConfirmationII}.

\begin{hbtheorem}{thm:p15-hempel-core-standards}{proved=annex:p15-hempel-core-standards}{The standards in the stated fragment}
For every admissible report in \cref{def:p15-hempel-essential}, the
satisfaction construction satisfies entailment, general consequence,
joint consistency and invariance under equivalent reports and hypotheses.
Consequently its confirmed, disconfirmed and neutral hypothesis classes
are mutually exclusive and exhaustive.
\end{hbtheorem}

The finite-domain model in the annex establishes consistency of the entire
confirmed class, not merely each displayed example. This proof uses the
absence of identity and unrestricted relational interpretations. The 1945
essay refers its more technical proofs to Hempel's 1943 article and explicitly
omits details. No claim is made here to have reconstructed that complete
technical article, resolved empty-essential-support conventions, or proved
an adequate scheme for equality, functions, higher-order variables or
quantified observation reports. Hempel himself lists richer languages and
quantified reports as further problems
\parencite[II, pp.~108--111, notes; pp.~120--121]{Hempel1945ConfirmationII}.

\begin{hbtheorem}{thm:p15-hempel-outside-name}{proved=annex:p15-hempel-outside-name}{Direct support and a consequence about an unmentioned name}
For $E=P(a)\land Q(b)$, the hypotheses $\forall x(P(x)\lor Q(x))$ and
$\exists xP(x)$ are directly confirmed. For a fresh name $c$, the sentence
$P(c)\lor Q(c)$ is confirmed, but not directly confirmed. Thus general
confirmation is not just the finite development test applied to the final
hypothesis.
\end{hbtheorem}

\begin{hbtheorem}{thm:p15-hempel-likes}{proved=annex:p15-hempel-likes}{Positive and negative confirmation without a conclusive inference}
Let $H=\forall x\exists yL(x,y)$. The report $L(a,b)\land L(b,a)$
directly confirms $H$ without entailing it over arbitrary nonempty domains.
The report $\neg L(a,a)\land\neg L(a,b)$ directly confirms $\neg H$
without entailing $\neg H$. The latter is nonconclusive disconfirmation,
not a counterexample to $H$ in every possible larger domain.
\end{hbtheorem}

\begin{hbtheorem}{thm:p15-hempel-evidence-extension}{proved=annex:p15-hempel-evidence-extension}{The whole record can reverse the classification}
$P(a)$ confirms $\forall xP(x)$ in the stated construction.
The enlarged consistent report $P(a)\land\neg P(b)$ disconfirms it and
does not confirm it. Confirmation is therefore not preserved by arbitrary
addition of observation sentences. Confirmation of a weaker hypothesis
from the \emph{same} report is a different condition.
\end{hbtheorem}

\section{Verification, acceptance and the limits of the proposal}
\label{sec:p15-hempel-verification}

\begin{hbdefinition}{def:p15-hempel-verification}{Relative verification and falsification}
$E$ relatively verifies $H$ when $E\models H$ and relatively falsifies it
when $E\models\neg H$. These are conclusive cases. They must not be
identified with definitive scientific acceptance or irrevocable rejection
of a report which might later be withdrawn.
\end{hbdefinition}

\begin{hbtheorem}{thm:p15-hempel-verification-cases}{proved=annex:p15-hempel-verification-cases}{Three contingent cases on infinite unrestricted domains}
For consistent finite ground reports in the unrestricted relational
language, $\exists xP(x)$ has a verifying report but no falsifying report;
$\forall xP(x)$ has a falsifying report but no verifying report; and
$\forall x\exists yL(x,y)$ has neither kind of report. These conclusions
concern these specified contingent sentences, not their displayed
quantifier prefixes alone.
\end{hbtheorem}

Hempel qualifies his general classification: a mixed-prefix sentence such
as $\forall x\exists y(P(x)\lor Q(y))$ is verified by $Q(a)$, while
$\forall x\exists y(P(x)\land Q(y))$ is falsified by $\neg P(a)$.
Analytic or contradictory matrices also cannot be classified merely by
their quantifier spelling
\parencite[II, \S10, pp.~112--114 and p.~113, n.~1]{Hempel1945ConfirmationII}.

For absolute verification Hempel makes the dependence on the status of
observation reports explicit. If none is accepted irrevocably, a conclusive
logical relation to a currently accepted report does not promise irrevocable
acceptance. If some are stipulated irrevocable, their entailed or
incompatible hypotheses inherit the corresponding definitive status.
He presents this as a choice in a reconstruction of scientific practice,
not a theorem inferred from a survey of actual laboratory behaviour.
The chess-rule analogy permits occasional violations of reconstructed rules
without identifying a rule with every observed move
\parencite[II, pp.~114--119]{Hempel1945ConfirmationII}.

In the Russell discussion, Hempel explicitly retracts the identification
of truth with high confirmation made in his earlier writings. Semantic
truth and evidential support can differ: either a hypothesis or its denial
can be true while neither is well confirmed by the accepted record.
Experimentally obtaining a report, logically comparing it with a hypothesis,
and deciding whether to accept the hypothesis remain distinct. His comments
on relative versus absolute falsifiability and on the restriction to
falsifiable hypotheses belong to this 1945 argument; they are not a
replacement for Popper's own presentation
\parencite[II, pp.~115--116, n.~2; pp.~119--121]{Hempel1945ConfirmationII}.

\begin{proofboundary}{def:p15-hempel-boundary}{What this method does and does not settle}
The task requires an explicit observational language, a complete stated
report and a hypothesis. Development supplies a finite named-class test;
consequence closure changes what can follow from direct confirmation.
The derivations above establish those specified logical properties and
separating cases. They establish neither a probability measure, a uniquely
projectible vocabulary, a reliable experimental record nor a generally
adequate acceptance policy. Source-specific quantitative accounts,
Hempel's complete 1943 technical construction and Sprenger's full 2011
historical argument are distinct presentations, not conclusions of the
nonempty named-evidence construction above.
\end{proofboundary}

\index{confirmation!Hempel}\index{raven paradox}\index{Nicod's criterion}
\index{development!of a hypothesis}\index{verification!relative and absolute}
